% ============================================================================
% Section 2: Related Work
% ============================================================================

\section{Related Work}
\label{sec:related}

Our work synthesizes techniques from five research areas: metamorphic
testing, mutation testing, statistical hypothesis testing, LLM/agent
evaluation, and probabilistic model checking. We survey each and
position \agentassay{} at their intersection.

\subsection{Software Testing Foundations}
\label{sec:related:testing}

\paragraph{Metamorphic Testing.}
Metamorphic testing~\citep{chen2018metamorphic,segura2016survey} addresses
the oracle problem by defining \emph{metamorphic relations} (MRs):
properties that must hold across related inputs. If $f(x_1)$ is unknown
but the relation $R(f(x_1), f(x_2))$ is known, a violation signals a
fault. Metamorphic testing has been applied to compilers, databases, and
search engines, and more recently to deep neural networks via
DeepTest~\citep{tian2018deeptest}. METAL~\citep{deng2024metal} extends
MRs to single-call LLM interactions but does not address multi-step
agent workflows where tool selection, reasoning chains, and
orchestration logic create a combinatorial space of execution paths.
\agentassay{} defines four families of agent-specific MRs that account
for the unique structure of agent traces (\cref{sec:metamorphic}).

\paragraph{Mutation Testing.}
Mutation testing~\citep{jia2011mutation,demillo1978hints} evaluates test
suite quality by injecting small syntactic faults (\emph{mutants}) into
the program and measuring how many the test suite detects
(\emph{kills}). The mutation score---the fraction of killed mutants---is
a measure of test adequacy. Traditional mutation operators target source
code (statement deletion, operator replacement, constant perturbation).
For AI agents, the ``source code'' includes prompts, tool
configurations, model selections, and context windows---none of which are
traditional code. \agentassay{} introduces four novel classes of mutation
operators tailored to these agent-specific artifacts (\cref{sec:mutation}).

\paragraph{Coverage Metrics.}
Code coverage~\citep{zhu1997software}---statement, branch, path, and
condition coverage---is the standard measure of testing thoroughness for
deterministic software. For agents, ``code'' is replaced by a
combination of prompts, tool inventories, and orchestration graphs, and
the execution space is stochastic rather than deterministic. No prior
work defines coverage metrics for agent execution. \agentassay{}
introduces a five-dimensional coverage tuple that captures tool
utilization, decision-path exploration, state-space coverage, boundary
testing, and model diversity (\cref{sec:coverage}).

\paragraph{The Oracle Problem.}
Barr et al.~\citep{barr2015oracle} survey the oracle problem: the
difficulty of determining correct output for a given input. For AI
agents, the oracle problem is acute---there is often no single ``correct''
output, only a distribution of acceptable behaviors. \agentassay{}
addresses this through three complementary oracle mechanisms: behavioral
contracts (\agentassert{}), metamorphic relations, and statistical
regression baselines.

\subsection{Statistical Testing and Sequential Analysis}
\label{sec:related:statistics}

\paragraph{Hypothesis Testing.}
The Neyman-Pearson framework~\citep{neyman1933problem,lehmann2005testing}
provides the mathematical foundation for deciding between two hypotheses
at controlled error rates. We adopt this framework for regression
detection: $H_0$ (no regression) vs.\ $H_1$ (regression occurred), with
significance level $\alpha$ controlling false alarms and power $1-\beta$
controlling missed regressions.

\paragraph{Sequential Analysis.}
Wald's Sequential Probability Ratio Test
(SPRT)~\citep{wald1947sequential} enables hypothesis testing with an
\emph{adaptive} sample size: testing stops as soon as sufficient evidence
accumulates, rather than requiring a fixed number of trials. This is
critical for agent testing, where each trial may cost \$0.01--\$1.00 in
API calls. We adapt SPRT for agent regression detection in
\cref{sec:stochastic:sprt}.

\paragraph{Confidence Intervals.}
The Clopper-Pearson interval~\citep{clopper1934use} provides exact
coverage for binomial proportions, while the Wilson
interval~\citep{wilson1927probable} offers better practical properties
for small samples. Agresti and Coull~\citep{agresti1998approximate}
demonstrate that approximate intervals often outperform exact ones. We
use the Wilson score interval as our primary confidence interval
construction, with Clopper-Pearson for formal guarantees
(\cref{sec:stochastic:verdict}).

\paragraph{Statistical Guidance for Software Engineering.}
Arcuri and Briand~\citep{arcuri2011practical} provide practical
guidelines for using statistical tests in software engineering
experiments, emphasizing effect sizes alongside $p$-values. We follow
their recommendations and require both statistical significance
\emph{and} practical significance (effect size $> \delta$) for regression
verdicts (\cref{sec:stochastic:regression}).

\subsection{LLM and Agent Evaluation}
\label{sec:related:evaluation}

\paragraph{LLM Evaluation Frameworks.}
deepeval~\citep{deepeval2024} provides metrics (faithfulness, answer
relevancy, hallucination rate) for RAG and conversational AI.
promptfoo~\citep{promptfoo2024} enables prompt-level testing with
assertion-based evaluation. OpenAI Evals provides a benchmark registry
for capability assessment. These tools evaluate \emph{output quality at
a point in time} but do not formalize \emph{regression detection across
versions}. They lack stochastic test semantics, coverage metrics, and
mutation testing.

\paragraph{Agent-Specific Testing.}
agentrial~\citep{agentrial2026} is the closest related work. It runs
agents multiple times, computes Wilson confidence intervals, applies
Fisher's exact test for regression detection, and includes CUSUM drift
detection. It supports seven framework adapters and CI/CD integration.
However, agentrial has no formal theoretical foundation: no stochastic
test semantics with provable guarantees, no coverage metrics, no
mutation testing, no metamorphic relations, no composition theory, no
SPRT for efficient testing, and no published paper. \agentassay{}
provides the formal foundation that agentrial lacks while also
introducing the six capabilities listed above. \Cref{tab:comparison}
provides a detailed feature comparison.

\begin{table}[t]
\centering
\caption{Feature comparison of agent testing approaches. \checkmark{}
indicates full support, $\circ$ partial support, --- no support.}
\label{tab:comparison}
\small
\begin{tabular}{lcccccc}
\toprule
\textbf{Feature} & \rotatebox{70}{\agentassay{}} &
  \rotatebox{70}{agentrial} &
  \rotatebox{70}{deepeval} &
  \rotatebox{70}{promptfoo} &
  \rotatebox{70}{METAL} &
  \rotatebox{70}{ProbTest} \\
\midrule
Multi-step agent workflows & \checkmark & \checkmark & --- & --- & --- & --- \\
Stochastic verdicts & \checkmark & $\circ$ & --- & --- & --- & $\circ$ \\
Formal test semantics & \checkmark & --- & --- & --- & --- & $\circ$ \\
Confidence intervals & \checkmark & \checkmark & --- & --- & --- & --- \\
Regression detection & \checkmark & \checkmark & --- & --- & --- & --- \\
SPRT adaptive stopping & \checkmark & --- & --- & --- & --- & --- \\
Coverage metrics & \checkmark & --- & --- & --- & --- & --- \\
Mutation testing & \checkmark & --- & --- & --- & --- & --- \\
Metamorphic relations & \checkmark & --- & --- & --- & \checkmark & --- \\
Contract integration & \checkmark & --- & --- & --- & --- & --- \\
Composition theory & \checkmark & --- & --- & --- & --- & --- \\
Bayesian analysis & \checkmark & --- & --- & --- & --- & --- \\
Published paper & \checkmark & --- & --- & --- & \checkmark & \checkmark \\
\bottomrule
\end{tabular}
\end{table}

\paragraph{LLM Repetition Studies.}
Raji et al.~\citep{raji2025repetitions} demonstrate empirically that
single-run LLM evaluations are unreliable and that multiple repetitions
are necessary to establish stable quality estimates. Their findings
motivate the multi-trial approach central to \agentassay{}, but they do
not formalize the statistical framework for determining how many
repetitions are sufficient or how to interpret the resulting
distribution.

\subsection{Agent Frameworks and Reliability}
\label{sec:related:frameworks}

The rapid growth of agent frameworks---AutoGen~\citep{wu2023autogen},
LangGraph~\citep{langgraph2024}, CrewAI~\citep{crewai2024}, OpenAI
Agents SDK~\citep{openai2025agents}---has created an ecosystem where
agents are composed, deployed, and updated at high velocity. None of
these frameworks include built-in regression testing capabilities.
Bhardwaj~\citep{bhardwaj2026abc} introduced Agent Behavioral Contracts
(ABC) for runtime enforcement of agent behavior via the \agentassert{}
framework, defining drift metrics and behavioral specifications.
\agentassay{} complements \agentassert{} by addressing the
\emph{pre-deployment} verification question: does the agent still satisfy
its contracts after a change? Separately, the growing reliance on
third-party agent skills and tool plugins has raised supply chain
security concerns~\citep{bhardwaj2026skillfortify}, reinforcing the
need for comprehensive quality assurance pipelines that span security
validation, behavioral testing, and runtime enforcement.

\subsection{Probabilistic Model Checking}
\label{sec:related:modelchecking}

PRISM~\citep{kwiatkowska2011prism} enables formal verification of
probabilistic systems through model checking of Markov decision processes
and continuous-time Markov chains. While PRISM provides strong
guarantees, it requires an explicit state-space model that is infeasible
to construct for LLM-based agents (the state space is effectively
infinite due to natural language). \agentassay{} takes a
\emph{testing-based} rather than \emph{verification-based} approach,
using statistical sampling to establish probabilistic guarantees without
requiring an explicit model.

\subsection{Flaky Tests}
\label{sec:related:flaky}

Research on flaky tests~\citep{luo2014empirical,parry2022survey}---tests
that non-deterministically pass or fail without code changes---reveals
that 2--16\% of test failures in large software projects are flaky.
Traditional approaches treat flakiness as a \emph{defect} to be
eliminated. \agentassay{} inverts this perspective: for agents,
non-determinism is a \emph{feature} of the system under test, and the
testing framework must accommodate it as a first-class concern rather
than attempting to suppress it.

\subsection{Positioning}
\label{sec:related:positioning}

\agentassay{} occupies a unique position at the intersection of formal
software testing theory and practical AI agent engineering. It is the
first framework to:
\begin{enumerate}[leftmargin=*]
  \item Provide formal stochastic test semantics with provable soundness
    and power guarantees for agent testing.
  \item Define coverage metrics, mutation operators, and metamorphic
    relations specifically designed for the agent execution model.
  \item Integrate with behavioral contracts for contract-as-oracle testing.
  \item Adapt sequential analysis (SPRT) for cost-efficient agent testing.
\end{enumerate}

No prior work---in software engineering, AI/ML evaluation, or formal
methods---addresses all of these concerns in a unified framework.
